
% ============================================================
%  3. Implémentation  (~3 pages)
% ============================================================
\section{Implémentation}

\subsection{Structure des packages}
Après avoir mis au point un ensemble l'ensemble des classes dans le diagramme de classe ci-dessus nous avons défini+.  l'architecture suivante : 
\dirtree{%
.1 src/.
.2 modele/.
.3 BaseDeDonnees.java
.2 expressions/.
.2 commandes/.
.2 parseur/.
.3 SGBD.jj
.2 execution/.
.2 persitance/.
.2 exceptions/.
}


\subsection{Package Modèle}
% Décrire les classes principales et leur rôle
% Mettre un extrait de code pertinent (ex: prochainSerial)

\begin{lstlisting}[caption=Compteur SERIAL global dans BaseDeDonnees]
public int prochainSerial() {
    this.compteurSerial++;
    return this.compteurSerial;
}
\end{lstlisting}

\subsection{Expressions logiques et évaluation WHERE}
% Expliquer le fonctionnement de ConditionSimple.evaluer()
% Mentionner le contexte ctx (Map<String, Tuple>)

\subsection{Analyse syntaxique avec JavaCC}
% Expliquer brièvement le rôle de JavaCC
% Montrer un extrait de la grammaire .jj

\begin{lstlisting}[language=bash, caption=Extrait de la grammaire SGBD.jj]
// Exemple de regle syntaxique
Commande parseCommande() :
{}
{
    ( parseCreate() | parseInsert() | parseSelect()
    | parseDrop()   | parseDelete() | parseTables()
    | parseExit() )
}
\end{lstlisting}

\subsection{Persistance}
% Expliquer le mécanisme de sauvegarde SQL Dump
% Montrer un exemple de fichier généré

\begin{lstlisting}[caption=Exemple de fichier de sauvegarde mabase.sql]
CREATE TABLE heroes (id SERIAL, nom VARCHAR, puissance INT);
INSERT INTO heroes (id, nom, puissance) VALUES (1, "Batman", 100);
INSERT INTO heroes (id, nom, puissance) VALUES (2, "Robin", 60);
\end{lstlisting}